\documentclass[10pt,a4paper]{amsart}
\usepackage[margin=19mm]{geometry}
\usepackage{amsmath,amssymb,mathtools}
\usepackage[scr=rsfs,cal=euler]{mathalfa}
\usepackage{xfrac}
\usepackage[unicode,hidelinks,bookmarks=false]{hyperref}
\newcommand{\ie}{i.e.\ }
\newcommand{\cf}{cf.\ }
\newcommand{\resp}{respectively\ }
\newcommand{\Subsection}{\subsection}
\providecommand{\nobreakdash}{\nobreak\hbox{-}}
\setlength{\emergencystretch}{2em}
\multlinegap0pt
\usepackage{bm}
\usepackage{bbm}
\usepackage{dsfont}
\usepackage{xspace}
\usepackage{cancel}
\usepackage{mathtools}
\usepackage{amsthm}
\makeatletter
\@ifundefined{theorem}{\newtheorem{theorem}{Theorem}}{}
\makeatother
\newcommand{\cg}{{\rm C}_{\geqslant}}
\newcommand{\mc}{\overline{\chi}_{\geqslant}}
\title[Majority C-coloring of graphs]{Majority C-coloring of graphs}
\author{Csilla Bujtas, Magda Dettlaff, Hanna Furmanczyk, Aleksandra Laskowska}
\date{Source: arXiv:2604.20752v1 (CC BY 4.0)}
\begin{document}
\maketitle

\noindent\emph{Source and licence.} Majority C-coloring of graphs. Version arXiv:2604.20752v1 (CC BY 4.0), distributed under the Creative Commons Attribution 4.0 International licence, as recorded in the arXiv licence metadata for this exact version and shown on the abstract page https://arxiv.org/abs/2604.20752v1. Full rights notice: licenses/arxiv-2604.20752-CC-BY-4.0.txt.\par\smallskip

\noindent\textit{Editorial scope.} Counted unit: the Theorem whose source label is \texttt{thm:M(n,k)} in the article (source0.tex lines 831--840, source label \texttt{thm:M(n,k)}), delivered together with the article's own complete original proof of it (source0.tex lines 841--894, one \texttt{proof} environment of 54 lines). No mathematics is written, repaired or re-derived: statement and proof are the article's own text line for line. Required context retained uncounted: eq:majority-1 (lines 303--305); eq:majority-2 (lines 309--311); thm:M(n,2) (lines 776--785). Every definition, display and label of those lines is retained unchanged. Omitted from this package and not redistributed: 1--302, 306--308, 312--775, 786--830, 895--1527. Nothing inside a retained excerpt is omitted or altered: every retained body is delivered exactly as the article's lines are written, so no omission receipt of any kind accompanies this package. Citations: the retained excerpts carry no citation command at all, so no bibliography entry of the article is transcribed and no reference is redistributed. Reference adaptation: none was needed and none was made; every reference inside the delivered unit resolves inside it, so no unresolved reference or citation is produced. Printed slips kept literal: no printed word, symbol, hypothesis or number of the counted statement or of its proof was altered. Layout and class: the article's own class is \texttt{article} with packages including graphicx, ulem, amsmath, amssymb, amsthm, soul, pgf, tikz, calc, enumitem, hyperref, verbatim, geometry, xcolor; the wrapper is the lane's pinned \texttt{amsart} editorial wrapper at 10pt on A4, which restates the theorem-like environments the retained text needs and adds the article's own macro definitions that the retained text needs (\texttt{\textbackslash cg}, \texttt{\textbackslash mc}), with extra wrapper packages (bm, bbm, dsfont, xspace, cancel, mathtools). The delivered numbering is the unit's own. Assets: no retained span contains a figure environment, a graphics-inclusion command, a PDF-page inclusion, a table or a TikZ picture; no image file is copied into this package, the package declares no asset dependency and no image file is redistributed with it. Licence: version arXiv:2604.20752v1 of this article is distributed under the Creative Commons Attribution 4.0 International licence, as recorded in the arXiv licence metadata for this exact version and shown on the abstract page https://arxiv.org/abs/2604.20752v1. A full rights notice is delivered at licenses/arxiv-2604.20752-CC-BY-4.0.txt. Characters: every retained excerpt is pure ASCII, so the delivered engine, XeLaTeX with its default Latin Modern text font, reports no missing character.
\begin{equation} \label{eq:majority-1}
    \bigl|\{\, u \in N(v) : \phi(u) = \phi(v) \,\}\bigr| \ge \frac{\deg_G(v)}{2}.
\end{equation} 

\begin{equation} \label{eq:majority-2}
    |N(v) \cap V_i| \ge |N(v) \cap (V\setminus V_i)|.
\end{equation}

   \begin{theorem} \label{thm:M(n,2)}
       For every integer $n \ge 2$, the maximum number of edges in an $n$-vertex graph $G$ with $\mc(G) =2$ is 
       $$
    M(n,2)= 
    \begin{cases}
	\frac{n(n-2)}{2}, &\text{if $n$ is even},\\
	\frac{(n-1)(n-2)}{2}, &\text{if $n$ is odd}.\\
    \end{cases}
    $$
   \end{theorem}

   \begin{theorem} \label{thm:M(n,k)}
       For every two integers $n \ge k \ge 2$, the maximum number of edges in an $n$-vertex graph $G$ with $\mc(G) =k$ is 
      \begin{equation} \label{eq:thm-M(n,k)}
    M(n,k)= 
    \begin{cases}
	{n-k+1 \choose 2 }, &\text{if $n-k$ is odd},\\
	{n-k+2 \choose 2 }- \frac{n-k+2}{2}, &\text{if $n-k$ is even}.\\
    \end{cases}
    \end{equation}
   \end{theorem}

   \begin{proof}
   First remark that, for $k=2$, these formulas are equivalent to those in Theorem~\ref{thm:M(n,2)}. Hence the statement holds for $k=2$ and any $n \ge 2$. We proceed by induction on $k$ by assuming  $k\ge 3$ and that the statement is true for the smaller values of $k$. 
   
   Let $G$ be a graph with $n$ vertices and $\mc(G) \ge k$. Let $V_1, \dots V_k$ be the (nonempty) color classes in a $\cg$-coloring of $G$ and $n_i=|V_i|$ for $i \in [k]$. We assume, without loss of generality, that $n_1 \leq n_2 \leq \cdots \leq n_k$. By $H(x)$, we denote the maximum number of edges in an $n$-vertex graph $G$ when $n_1=x$ and $\mc(G) \ge k$. Observing that $x \in \{1, \dots , \lfloor \frac{n}{k} \rfloor \}$, we distinguish the following three cases. 
   \begin{description}
       \item[Case 1. $n_1=1$] 
       By the majority condition~\eqref{eq:majority-1}, the vertex $v_1$ in $V_1$ is an isolated vertex of $G$. Therefore, $\mc(G) \ge k$ implies $\mc(G-v_1) \ge k-1$ and we can apply the hypothesis for $G-v_1$. If $n-k$ is odd, then so is $(n-1)-(k-1)$ and we obtain
       $$H(1) \leq {(n-1)- (k-1) +1 \choose 2}= {n-k+1 \choose 2}.
       $$
       If $n-k$ is even, then so is $(n-1)-(k-1)$ and therefore,
       \begin{align*}
           H(1) &\leq {(n-1)- (k-1) +2 \choose 2}- \frac{(n-1)-(k-1)+2}{2}\\
       &= {n-k+2 \choose 2 }- \frac{n-k+2}{2}.
       \end{align*}
         In both cases, the formulas in~\eqref{eq:thm-M(n,k)} are upper bounds on $H(1)$.  
       \item[Case 2. $2 \leq n_1 <\frac{n}{k}$]  
       We modify $G$ to obtain a graph $G'$ with a $\cg$-coloring that uses $k$ colors and the smallest color class contains only one vertex. We also ensure that the size of $G$ does not decrease in the process.

       Let $V_1, \dots, V_k$ be the color classes for $G$. Choose an arbitrary vertex $v_1$ from $V_1$ and let $X=V_1 \setminus \{v_1\}$. Since $n_1 \ge 2$, the cardinality of $X$, denoted by $x$, is at least $1$. To construct $G'$, delete all edges incident to $v_1$ and add all edges (that were missing in $G$) between $X$ and $V_k$. 

       We first prove that $m(G)$ and $m(G')$, the number of edges in $G$ and $G'$ respectively, satisfy $m(G) \leq m(G')$. Vertex $v_1 $ has at most $n_1-1$ neighbors in $V_1$. By the majority condition~\eqref{eq:majority-1}, $\deg_G(v_1) \leq 2n_1-2$ follows. Therefore, we deleted at most $2n_1-2$ edges from $G$. By the same majority condition, every vertex $u \in X$ has at most $n_1-1$ neighbors from $V_k$ in $G$. The condition $n_1 < \frac{n}{k}$ implies $n_k \ge n_1+1$, and consequently, at least two new edges incident to $u$ were added to obtain $G'$. Therefore, the number of new edges in $G'$ is at least $2x$, and we may derive 
       $$m(G') \ge m(G) -(2n_1-2) + 2x = m(G). 
       $$
       Next, we show that $V_1', \dots, V_k'$, where $V_1'=\{v_1\}$, $V_k'=V_k \cup X$, and $V_i'=V_i$ for $i\in \{2, \dots, k-1\}$, defines a $\cg$-coloring for $G'$. The majority condition~\eqref{eq:majority-2} clearly holds for the isolated vertex $v_1$. For the vertices in $\bigcup_{i=2}^{k-1} V_i$, the number of neighbors from the same color class remained unchanged, and the neighbors from different classes remained the same or decreased by $1$, and~\eqref{eq:majority-2} still holds in $G'$. For every vertex $u \in X$, since the majority condition holds for $u$ in $G$, and further by the construction of $G'$, we have that
       \begin{align*}
           |N_{G'}(u) \cap V_k'| &\ge |N_{G}(u) \cap V_1 | \ge \Bigl|N_{G}(u) \cap \bigcup_{i=2}^{k-1} V_i \Bigr|\\
           &= \Bigl|N_{G'}(u) \cap \bigcup_{i=2}^{k-1} V_i' \Bigr|
       = \Bigl|N_{G'}(u) \cap \bigcup_{i=1}^{k-1} V_i' \Bigr|.
       \end{align*}
       That is, $u$ satisfies the condition~\eqref{eq:majority-2} in $G'$. The last case to check is when $u \in V_k$. During the construction of $G'$, vertex $u$ might gain some new neighbors from $X \subseteq V_k'$, and it might lose one neighbor, namely $v_1$, which is outside $V_k'$. It implies that~\eqref{eq:majority-2} remains true for $u$ in $G'$. We conclude that $\mc(G') \ge k$ and there is a $\cg$-coloring of $G'$ where the smallest color class consists of an isolated vertex.

       Consequently, for every graph $G$ satisfying the conditions of this case, we can construct a graph $G'$ with $m(G') \ge m(G)$ such that a $\cg$-coloring of $G'$ belongs to Case~1. Formally,
       $$m(G) \leq m(G') \leq H(1)$$
       and hence, the relevant formula in~\eqref{eq:thm-M(n,k)} gives an upper bound on the size of $G$.
       \item[Case 3. $n_1=\frac{n}{k}$]
       In this case, the color classes are of the same cardinality $\frac{n}{k}$. Given a graph $G$ that admits such a $\cg$-coloring, we construct $G'$ as in Case~2. The only difference is that instead of always choosing the color class $V_k$ to be replaced with $V_k \cup X$, we choose a color class $V_s$ with $s \ge 2$ such that the number of edges between $V_s$ and $X$ is the smallest in $G$. 
       
       By the majority condition~\eqref{eq:majority-2}, there are at most $(n_1-1)^2$ edges between $X$ and $\bigcup_{i=2}^k V_i$ in $G$. Since $k \ge 3$, the number of edges between $X$ and $V_s$ is at most $\lfloor \frac{1}{2} (n_1-1)^2 \rfloor$. When constructing $G'$, the at most $2n_1-2$ edges incident to $v_1$ are removed, and at least
       $(n_1-1)|V_s|- \lfloor \frac{1}{2} (n_1-1)^2 \rfloor$ new edges are added between $X$ and $V_s$. If $n_1 \ge 3$, then the size of $G'$ can be estimated as 
       \begin{align} \label{eq:case-3}
                m(G') &\ge m(G) -(2n_1-2) + (n_1-1)n_1 - \Bigl\lfloor \frac{(n_1-1)^2}{2} \Bigr\rfloor\\
        \nonumber  &\ge m(G) + \frac{n_1^2}{2}-2n_1+\frac{3}{2} \ge m(G).
         %&=  m(G) + \frac{1}{2}(n_1-2)^2-0.5          
       \end{align} 
       If $n_1=2$, then $\lfloor \frac{1}{2} (n_1-1)^2 \rfloor=0$ and a direct substitution into~\eqref{eq:case-3} shows $m(G') \ge m(G)$. 

       What remains to check is that the partition into $V_1'=\{v_1\}$, $V_s'= V_s \cup X$, and $V_i'=V_i$ for $i \in [k] \setminus \{1,s\}$ defines a $\cg$-coloring of $G'$. It can be done along the same lines as for Case~2.
       We may infer again that $H(1)$, and consequently also the relevant formula from~\eqref{eq:thm-M(n,k)}, give an upper bound on $m(G)$.
   \end{description}
   The three cases above cover all possibilities. We therefore conclude that~\eqref{eq:thm-M(n,k)} gives an upper bound on the number of edges for every graph $G$ with  $\mc(G) \ge k$. 
   \medskip
   
   To complete the proof, we show that the upper bound is attained by a graph $G_{n,k}$ for every $n \ge k$ such that $\mc(G_{n,k})=k$ also holds. If $n-k$ is odd, let $G_{n,k}$ be the graph consisting of $k-1$ isolated vertices and a component which is a complete graph on $n-k+1$ vertices.  If $n-k$ is even, we take $k-2$ isolated vertices and a complete graph of order $n-k+2$ from which a perfect matching is removed to obtain $G_{n,k}$. It is straightforward to check that $\mc(G_{n,k})=k$ holds in both cases and the size of  $G_{n,k}$ corresponds to the values in~\eqref{eq:thm-M(n,k)}. 
   \end{proof}

\end{document}
